%!TeX root=MemoriaTFG.tex

\chapter{Material complementario y repositorio reproducible}
\label{cap:repo}

El cuerpo de esta memoria se circunscribe a la pregunta de
investigación central: el rendimiento agregado y la capacidad de
generalización de los modelos fundacionales para imagen
dermatológica. El material instrumental, las tablas extendidas de
resultados, los análisis complementarios y los artefactos de
software derivados del trabajo se publican de forma estructurada en
un repositorio reproducible público, con el fin de no fragmentar la
argumentación del documento principal y de habilitar la
verificación independiente de las cifras reportadas en el
capítulo~\ref{cap:resultados}.

\section{Repositorio público}
\label{sec:repo-general}

El conjunto del material complementario se distribuye en el
repositorio \texttt{dermapixel-tfg-EPS0270}, accesible a través de
la página de documentación

\begin{center}
\url{https://tcontesti.github.io/dermapixel-tfg-EPS0270/}
\end{center}

\noindent que actúa como portada e índice navegable del contenido.
El repositorio contiene el código de los experimentos, el
\emph{datasheet} del dataset DermapixelAI~1.0, las tablas detalladas
de resultados, la descripción del prototipo, las ablaciones
complementarias, el material de interpretabilidad y la declaración
de uso de herramientas de inteligencia artificial. Todas las cifras
publicadas en el repositorio coinciden de forma exacta con las
reportadas en el cuerpo de la memoria; no se introduce ningún
resultado adicional que altere las conclusiones del documento
principal.

\section{Estructura de carpetas}
\label{sec:repo-estructura}

El repositorio se organiza en directorios de responsabilidad
disjunta. La tabla~\ref{tab:repo-estructura} sintetiza su
contenido.

\begin{table}[H]
\centering
\footnotesize
\setlength{\tabcolsep}{5pt}
\renewcommand{\arraystretch}{1.05}
\caption{Estructura de directorios del repositorio reproducible.}
\label{tab:repo-estructura}
\begin{tabular}{ll}
\hline
Directorio & Contenido \\
\hline
\texttt{code/} & Pipelines de \emph{linear probing}, \emph{fine-tuning}, segmentación con LoRA y evaluación \emph{zero-shot}. \\
\texttt{datasheet/} & \emph{Datasheet} y documento de entrega del dataset DermapixelAI~1.0. \\
\texttt{tables/} & Tablas detalladas de resultados por dataset (equidad, evaluación cruzada, segmentación). \\
\texttt{ablations/} & Ablaciones complementarias sobre DermapixelAI~1.0 y rama contrastiva SpanDerm. \\
\texttt{sae/} & Diccionario de conceptos y modelos de interpretabilidad (\emph{Sparse Autoencoders}, CBM). \\
\texttt{llm/} & Comparación extendida de modelos de lenguaje multimodales. \\
\texttt{prototype/} & Descripción del prototipo DermApIxel y de sus componentes. \\
\texttt{repro/} & Mapa de tareas experimentales, recetas de entrenamiento, semillas y \emph{checkpoints}. \\
\texttt{ai-statement/} & Declaración sobre el uso de herramientas de inteligencia artificial. \\
\hline
\end{tabular}
\end{table}

\section{Índice de correspondencia}
\label{sec:repo-indice}

La tabla~\ref{tab:repo-indice} establece la correspondencia entre el
material complementario y su ubicación en el repositorio. Cada
entrada puede consultarse a través de la página de documentación
indicada en la sección~\ref{sec:repo-general}.

\begin{table}[H]
\centering
\footnotesize
\setlength{\tabcolsep}{5pt}
\renewcommand{\arraystretch}{1.1}
\caption{Índice de correspondencia entre el material complementario y el repositorio.}
\label{tab:repo-indice}
\begin{tabular}{lll}
\hline
Material complementario & Ubicación en el repositorio & Sección \\
\hline
Mapa de tareas y reproducibilidad & \texttt{repro/} & \ref{sec:repo-reproducibilidad} \\
Tablas detalladas de resultados & \texttt{tables/} & \ref{sec:repo-tablas} \\
Intervalos de confianza \emph{bootstrap} & \texttt{tables/} & \ref{sec:repo-bootstrap-ci} \\
Calibración \emph{endpoint} melanoma HAM10000 & \texttt{tables/} & \ref{sec:repo-calibracion-ham} \\
Pipeline y caracterización del dataset & \texttt{datasheet/}, \texttt{ablations/} & \ref{sec:repo-ablaciones} \\
Declaración de uso de IA & \texttt{ai-statement/} & \ref{sec:repo-ia} \\
Documento de entrega del dataset & \texttt{datasheet/} & \ref{sec:repo-datasheet} \\
\emph{Sparse Autoencoders} y conceptos & \texttt{sae/} & \ref{sec:repo-sae} \\
Modelos de lenguaje multimodales & \texttt{llm/} & \ref{sec:repo-llm} \\
Prototipo DermApIxel & \texttt{prototype/} & \ref{sec:repo-prototipo} \\
Ablaciones complementarias & \texttt{ablations/} & \ref{sec:repo-ablaciones} \\
\hline
\end{tabular}
\end{table}

\section{Datasheet del dataset DermapixelAI 1.0}
\label{sec:repo-datasheet}

El repositorio incluye el \emph{datasheet} completo del dataset
DermapixelAI~1.0 y el documento formal de entrega, con la
identidad, la procedencia, la licencia CC-BY-NC-SA~4.0, las
condiciones de disponibilidad, la composición, las
recomendaciones de uso responsable y el plan de mantenimiento del
recurso. La caracterización resumida del dataset figura en el
capítulo~\ref{cap:metodos}.

\section{Tablas detalladas de resultados}
\label{sec:repo-tablas}

El repositorio recoge las tablas extendidas de resultados que
desagregan, por dataset y por modelo, las cifras sintetizadas en el
capítulo~\ref{cap:resultados}: el desglose de equidad por fototipo,
la matriz de evaluación cruzada entre conjuntos, los resultados de
segmentación por arquitectura y las tareas de transferencia
adicionales.

\section{Intervalos de confianza por \emph{bootstrap}}
\label{sec:repo-bootstrap-ci}

El repositorio publica los intervalos de confianza al $95\%$ por
\emph{bootstrap} ($B = 1\,000$, remuestreo estratificado por clase,
percentil $2{,}5/97{,}5$, semilla $42$), celda a celda, de las
métricas \emph{headline} que en el cuerpo se reportan como
estimaciones puntuales: el sondeo lineal de PanDerm sobre los
conjuntos externos, la comparativa de paradigmas de razonamiento
clínico, el desglose de equidad por fototipo y el \emph{benchmark}
frente a líneas base convolucionales. El material documenta el
mecanismo de cómputo ---reajuste determinista de un clasificador
lineal sobre los \emph{embeddings} congelados, o remuestreo directo
sobre las predicciones guardadas, según el caso--- y declara
abiertamente las celdas no cubiertas y las convenciones de cálculo
empleadas por fila.

\section{Calibración en el \emph{endpoint} de melanoma de HAM10000}
\label{sec:repo-calibracion-ham}

El repositorio publica el detalle de la calibración de los tres
codificadores en el \emph{endpoint} binario melanoma frente al resto
sobre el test de HAM10000 (sección~\ref{sec:equidad-delong},
tabla~\ref{tab:delong-ham}): el error de calibración esperado
(ECE, $15$ \emph{bins} de igual anchura sobre la probabilidad de la
clase positiva) y el \emph{Brier score}, junto con las curvas de
fiabilidad por \emph{bin} que los sustentan. Las probabilidades por
muestra son exactamente las del test de DeLong de ese \emph{endpoint}
---reajuste determinista de la misma cabeza lineal sobre los
\emph{embeddings} congelados, sin reentrenamiento---; el cómputo
verifica previamente, celda a celda, que la AUROC recomputada reproduce
el valor publicado antes de admitir cualquier cifra de calibración.

\section{Ablaciones complementarias}
\label{sec:repo-ablaciones}

El repositorio documenta las ablaciones complementarias sobre el
dataset DermapixelAI~1.0, el pipeline de construcción y
caracterización del dataset, la ablación del modelo fundacional
inglés con traducción ES--EN y las iteraciones de la rama
contrastiva en castellano, incluyendo configuración de
entrenamiento, barridos de épocas, techos de sondeo lineal y
calibración.

\section{Sparse Autoencoders y diccionario de conceptos}
\label{sec:repo-sae}

El repositorio contiene el material de interpretabilidad: la
configuración de los \emph{Sparse Autoencoders} sobre el modelo
fundacional, el diccionario de conceptos clínicos, los modelos de
cuello de botella conceptual (\emph{Concept Bottleneck Models}) y
los coeficientes derivados de las anotaciones \emph{Seven-Point
Checklist}.

\section{Modelos de lenguaje multimodales}
\label{sec:repo-llm}

El repositorio presenta la comparación extendida de modelos de
lenguaje multimodales evaluados como sistemas de razonamiento
clínico, con el detalle de los \emph{prompts}, los proveedores y las
métricas que sintetiza el capítulo~\ref{cap:resultados}.

\section{Prototipo DermApIxel}
\label{sec:repo-prototipo}

El repositorio describe el prototipo DermApIxel, que integra los
componentes derivados del trabajo en un entorno operativo. Su
evaluación operativa queda fuera del alcance del presente documento;
el repositorio recoge la descripción de sus módulos, los resultados
de \emph{ensemble} y los análisis de recuperación visual y
codificadores adicionales.

\section{Declaración de uso de inteligencia artificial}
\label{sec:repo-ia}

El repositorio incluye la declaración sobre el uso de herramientas
de inteligencia artificial durante la elaboración del trabajo, con
el alcance, las herramientas empleadas y las responsabilidades del
autor.

\section{Reproducibilidad}
\label{sec:repo-reproducibilidad}

El repositorio documenta el mapa de tareas experimentales, las
recetas de entrenamiento, las semillas, las normalizaciones y los
\emph{checkpoints} necesarios para reproducir de forma independiente
las cifras del capítulo~\ref{cap:resultados} a partir del código
publicado.
